\section{Finite affine construction certificates}
\label{app:template-certificates}

Some construction families use a fixed edge trade whose vertices move
affinely with the arm lengths.  This section specifies the proof data,
the verification argument, and the exact parameter coverage of these
families.  In particular, fitting a formula to several examples is not
an existence argument here.  Every formula used in a construction is
checked against its source word on an entire integer ray.

The finite inputs are included in the accompanying machine-readable
supplement, with the names and hashes in Section~\ref{templates:ledger}.
All $16$ minimum-two templates and all $216$ minimum-four templates
are included in full.  The two equal-even repairs, the alternative
endpoint repair, and the four equal-odd four-arm repairs are also
printed in Section~\ref{templates:small-tables}, including every
changed vertex's source arm and depth.  The tables are construction
instructions; a search program is not needed to instantiate them.

\subsection{Finite word inventories}
\label{templates:finite-check}

A finite spider certificate specifies a hub label $z$, an ordered list
of arm lengths, and one complete outward word for each arm.  Put
$N$ equal to the sum of those lengths.  To check an ordinary graceful
certificate, check the word lengths and the two multiset identities
\[
 \{z\}\uplus\biguplus_i\{P_{i,j}\}=[0,N],\qquad
 \biguplus_{i,j}|P_{i,j}-P_{i,j-1}|=[1,N],\quad P_{i,0}=z.
\]
For a hub-zero certificate also require $z=0$.  For an offset-$k$
certificate require $z=0$ and replace the first identity on nonroot
labels by $\biguplus_i\{P_{i,j}\}=[k+1,k+N]$; the color identity
is unchanged.  For alpha, check $z=0$, the cut
$K=\sum_i\floor{L_i/2}$, and that depth $j$ is lower exactly when
$j$ is even.  A graceful permutation on $n$ vertices is checked by
labels $[0,n-1]$ and successive differences $[1,n-1]$, with the
stated first value and alpha condition when applicable.

All checks are finite list operations with integer arithmetic.  A
catalog must additionally have distinct profile keys and exactly
the required profile set.  Controlled interfaces, such as a pinned
short word or a specified upper-leaf distance, are checked directly
on the words.  Thus a finite record's search status or node count
has no role in its acceptance as a labeling proof.

\subsection{Affine data and source coordinates}
\label{templates:encoding}

An integer ray has the form $h\in\mathbb Z$, $h\ge h_0$.  A vertex
name is an affine function $Ah+B$, not a numerical vertex sampled at
a particular $h$.  The templates use the following three encodings.
\begin{enumerate}
\item In \texttt{S2o.jsonl}, \texttt{S2e.jsonl}, and
\texttt{S4.jsonl}, a pair \texttt{[A,B]} means $Ah+B$.
The first two files use $b=8h+r$; the third uses $b=48h+r$.
\item In \texttt{E.json} and \texttt{Alt.json}, a pair
\texttt{[A,B]} means $(Ar+B)/2$.  The parameter is
$r=4q+\epsilon$, $q\ge3$, with $\epsilon=1$ or $3$ as specified.
After substitution, every label and depth is integral affine in $q$.
\item In \texttt{Q0.json},\ldots,\texttt{Q3.json}, the nested pair
\texttt{[[a,b],[c,d]]} means $(a/b)q+c/d$.  Here $s=4q+r$,
$r\in\{0,1,2,3\}$, and $q\ge6$.  The rational representation is
exact; all actual label and depth coefficients become integers.
\end{enumerate}

Arm numbers are zero based.  Depth zero is the hub, and the edge of
depth $j$ joins depths $j-1$ and $j$.  In position tables the arm
number $-1$ denotes the hub, whose depth and label are zero.
For an explicitly defined source spider, write
$F_i(h,j)$ for the label at depth $j$ on source arm $i$, setting
$F_i(h,0)=0$.  The source formulas, including all exceptional small
path bases, are specified below and in the common construction
appendices.

\begin{center}
\begin{tabular}{p{.28\textwidth}p{.61\textwidth}}
Field & Mathematical meaning\\\hline
\texttt{residue}, \texttt{delta} & $r$ and the target long-arm gap $\delta$\\
\texttt{remove} & List of objects with fields \texttt{u}, \texttt{v},
\texttt{arm}, \texttt{depth}; these specify a source edge and its two affine labels\\
\texttt{add} & List of affine endpoint pairs specifying the replacement edges\\
\texttt{source\_b} & Parameter used during discovery; it is not the lower bound of the proved ray\\
\texttt{verification.h\_min} & The inclusive verified ray bound $h_0$\\
\texttt{verification.b\_min} & $8h_0+r$ or $48h_0+r$, respectively\\
\texttt{removed\_edges}, \texttt{added\_edges} & Remove/add lists in the equal-source and four-odd formats\\
\texttt{positions} & Affine label-to-source-coordinate table; a dictionary in E/Alt, an object list in the Q files\\
\texttt{q\_minimum} & Inclusive lower bound six in a Q recipe.
\end{tabular}
\end{center}
The cached \texttt{verification} output is descriptive.  A verifier
recomputes its obligations from the source words and the edge lists.
It does not accept the recorded word ``verified'' as a proof premise.
For the minimum-arm files its remaining output fields are
\texttt{removed\_edges}, \texttt{symbolic\_constraints}, and
\texttt{arm\_lengths}, together with copies of the residue and gap.
An E recipe is selected by the top-level key \texttt{"1"} or
\texttt{"3"}; an Alt recipe has the same three fields directly.
In their position dictionaries an entry
\texttt{"a,b":[i,[c,d]]} means label $(ar+b)/2$ at source arm $i$
and depth $(cr+d)/2$.  A Q position is instead an object with fields
\texttt{label}, \texttt{arm}, and \texttt{depth}, decoded using the
rational-coefficient convention above.

\subsection{Why the finite verification proves an infinite family}
\label{templates:soundness-section}

The source need not yet have the desired color inventory.  Let its
signed color defect be
\[
 D_0=\mathcal C(G_0)-[1,N],
\]
where $[1,N]$ denotes the multiset containing each integer once and
$\mathcal C(G_0)$ counts edge differences with multiplicity.  In the
applications below, an interval proof gives $D_0$ as a finite signed
sum of affine color functions.  For the minimum-two and minimum-four
sources, $D_0=0$.

\begin{theorem}[Soundness of an affine trade certificate]
\label{templates:soundness}
Suppose a source graph on labels $0,\ldots,N(h)$ is an explicitly
labeled spider for every integer $h\ge h_0$, all its edges cross a
specified cut $K(h)$, and its signed color defect is the proved
finite expression $D_0(h)$.  Let $R$ and $A$ be finite affine edge
lists.  The following checks imply that replacing $R$ by $A$ gives
a root-zero alpha labeling of the stated target spider for every
$h\ge h_0$.
\begin{enumerate}
\item All participating labels and depths are integral and in range.
Every recorded vertex is proved to equal its source coordinate
$F_i(h,j)$.  Distinct recorded vertices never collide on the ray.
\item Every edge of $R$ occurs at its recorded source arm and depth.
The lists contain no repeated edges, and every edge of $A$ is absent
from the retained source graph.  Each added edge crosses $K(h)$.
\item The endpoint-incidence multisets of $R$ and $A$ agree.  With
each edge oriented from its lower label to its upper label, its
color is the positive affine difference of its endpoint functions.
As a signed multiset of affine functions,
\begin{equation}\label{templates:color-identity}
 D_0-\mathcal C(R)+\mathcal C(A)=0.
\end{equation}
\item The weighted graph obtained by contracting all untouched path
intervals has exactly the stated number of hub-to-leaf paths, uses
every contracted vertex once, and has the stated target arm lengths
as affine identities.
\end{enumerate}
\end{theorem}

\begin{proof}
Fix an arbitrary integer $h\ge h_0$.  The coordinate checks identify
actual distinct vertices and actual removable edges of the source.
The retained and added edges therefore define the prescribed graph
on exactly the original vertex-label set.  In particular, every label
$0,\ldots,N(h)$ still occurs once, and the hub is still zero.
All retained and added edges cross the same cut.  Evaluating the
formal multiset identity \eqref{templates:color-identity} gives the
complete color multiset $[1,N(h)]$, even if two affine color functions
happen to coincide at that parameter.  The endpoint-incidence check
also preserves the degree of every labeled vertex.

To justify the topology check, on each source arm cut exactly the
edges in $R$.  Every remaining component is an interval of that path.
Keep its boundary vertices, and replace a nontrivial interval by a
weighted edge whose weight is its exact number of original edges.
A singleton component remains a vertex.  Different interval interiors
are disjoint.  Reintroduce the edges in $A$ with weight one.  Expanding
the weighted edges recovers precisely the traded graph.  Thus a
traversal proving that the weighted graph consists of the required
hub paths proves the same assertion for the full graph, with each
arm length equal to the sum of its weights.  Exhausting the weighted
vertices excludes any additional component or cycle.  This proves the
target topology, and hence the root-zero alpha labeling.  Since $h$
was arbitrary, the result holds on the whole ray.
\end{proof}

\subsubsection{Exact operations on an integer ray}

All affine assertions reduce to finite rational arithmetic.  Equality
means equality of both coefficients.  For $Ah+B\ge c$ on
$h\ge h_0$, it suffices and is necessary that $A\ge0$ and
$Ah_0+B\ge c$.  A strict positive inequality is treated with $c=1$
after integrality has been established.  Distinctness is checked by
solving $Ah+B=0$: a collision exists precisely when its root is an
integer at least $h_0$, apart from the constant cases.

When a source word has a piecewise coordinate formula, the verifier
first selects the appropriate finite base, block index, and residue
of the depth.  It then proves the inequalities defining that branch
on the full ray.  Any remaining periodic test is resolved by its
finitely many residue classes.  Ordering cut depths may be proposed
using a convenient parameter, but every consecutive difference in
the proposed order is subsequently proved positive on the whole
ray.  The proposed order alone supplies no mathematical inference.

The threshold computation used for the minimum-two and minimum-four
tables can be described explicitly.  Start with $h_0=0$.  Each
required inequality $Ah+B\ge c$ with $A>0$ raises the bound to
\[
 h_0\leftarrow\max\left(h_0,\ceil{(c-B)/A}\right).
\]
A constant inequality must already hold, and a negative slope fails
the selected eventual branch.  The inclusive bounds printed below
are sufficient certified bounds, not claims of optimal construction
thresholds.

\subsubsection{Verification pseudocode}

The finite procedure underlying Theorem~\ref{templates:soundness}
is the following.  Formal colors are stored as coefficient pairs
with signed multiplicities.
\begin{verbatim}
decode the parameter domain, source words, and affine records
for every participating vertex:
    prove integrality, source arm/depth, and label identity
prove all required ranges and vertex distinctness on the ray
for every removed edge:
    prove its two depths are consecutive on one source arm
for every added edge:
    prove distinctness, absence in the retained source, and cut crossing
compare the removed and added endpoint-incidence counters
cancel seed_defect - removed_color_counter + added_color_counter
on every source arm:
    sort and certify all cut depths
    contract each untouched interval with its exact affine length
insert each added edge with weight one
traverse every hub branch, rejecting repeated or omitted vertices
compare the affine path-length multiset with the target lengths
accept only if every preceding exact check succeeds
\end{verbatim}
For example, a removed edge with larger depth $d$ separates an
untouched interval ending at $d-1$ from one starting at $d$.
The weight between retained depths $u\le v$ is $v-u$; a zero
weight contributes a singleton, not a new edge.  This convention
avoids counting vertices as edges in the topology calculation.

\subsection{The fixed-endpoint path words used by the templates}
\label{templates:path-words}

Let $P_e(n)$ be an alpha word on $\{0,\ldots,n-1\}$ beginning at
the lower endpoint $e$, for $e\in\{1,2,3\}$.  The finite bases and
prefixes specify these words completely.  For $e=1$ the prefixes and
bases are
\[
 Q_1=(1,2,0,3),\qquad
 \begin{array}{c|l}
 n& P_1(n)\\\hline
 3&(1,2,0)\\
 4&(1,2,0,3)\\
 6&(1,5,0,3,2,4)\\
 9&(1,8,0,6,2,7,4,5,3).
 \end{array}
\]
For $e=2,3$, input \texttt{P.json} gives
\[
 Q_2=(2,5,3,4,0,7,1,6),\qquad
 Q_3=(3,11,0,10,1,8,2,7,4,6,5,9).
\]
Its base sizes are respectively
\[
 \begin{array}{c|l}
 e=2&5,6,7,8,10,11,12,17,\\
 e=3&7,8,9,10,11,12,14,15,16,17,18,25.
 \end{array}
\]
The file's top-level keys are \texttt{"2"} and \texttt{"3"}; each
has fields \texttt{prefix} and \texttt{bases}, the latter a dictionary
from the decimal base size to its entire vertex word.  Every one
of these finite words is checked by the ordinary permutation and
alpha inventories.

Here is a closed coordinate rule for all larger words.  Choose the
unique base size $\beta$ congruent to $n$ modulo $4e$, and write
$n=\beta+4ez$.  For $j=0,\ldots,z-1$, replace a value $x$ of
$Q_e$ by
\begin{equation}\label{templates:prefix-coordinate}
 \begin{cases}
 x+2ej,&x<2e,\\
 x+n-4e-2ej,&x\ge2e.
 \end{cases}
\end{equation}
Concatenate these $z$ blocks, followed by $P_e(\beta)+2ez$.
Equivalently, one prefix extends a word of order $n$ to order
$n+4e$: raise the upper prefix labels by $n$, then append the old
word translated by $2e$.  The prefix colors are
$n+1,\ldots,n+4e-1$, the joining color is $n$, and the old colors
are $1,\ldots,n-1$.  The lower and upper vertex blocks are
consecutive and disjoint.  This proves the recurrence, including
its first endpoint.  The complete permitted domains are
$n\ge2e+1$, $n\ne4e+1$; the exceptional residue begins instead at
$n=8e+1$.

\subsection{Sixteen minimum-two templates}
\label{templates:min2}

Write $b=8h+r$, $0\le r<8$, and let $\delta$ be the long-arm gap.
For odd $r$, $\delta\in\{0,1\}$; for even $r$,
$\delta\in\{1,2\}$.  Thus there are exactly $8+8=16$ templates.
The even repeated case $\delta=0$ is proved separately by the
balanced-even construction in Section~\ref{app:small-alpha}.

To define the source, put $n=b+\delta+2$ and $k=\floor{b/2}$.
Start with
\[
 (b-1,1),\qquad (b,\,P_1(b-3)+2),
\]
where the notation in the second word means concatenation of $b$
with the translated path.  Raise its labels above $k$ by $n$.
The added $n$-vertex arm is
\begin{equation}\label{templates:min2-newword}
 \begin{cases}
 (b+2+j,b+1-j)_{j=0}^{b/2},&b\text{ even},\ \delta=0,\\
 \operatorname{rev}(P_1(n))+k+1,&n\text{ even},\\
 n+k-\operatorname{rev}(P_1(n)),&n\text{ odd}.
 \end{cases}
\end{equation}
The first line specifies the separately used even repeated source;
the $16$ repair records use the appropriate other line.  The source
arm lengths are $(2,b-2,b+\delta+2)$.  Its labels, colors, and
alpha separation follow from the path recurrence and the middle
interval attachment; its defect is zero.  Its total size is
$N=2b+\delta+2$ and the target cut is
$K=1+\floor{b/2}+\floor{(b+\delta)/2}$.

Each record's \texttt{remove} entry identifies an actual edge in
one of these three source arms.  Applying its \texttt{add} list and
the checks of Theorem~\ref{templates:soundness} gives arm lengths
$(2,b,b+\delta)$ as exact affine identities.  The files
\texttt{S2o.jsonl} and \texttt{S2e.jsonl} have respectively eight
records; their exact inclusive thresholds and edge counts are in
Section~\ref{templates:thresholds}.  For illustration, the first
minimum-two odd record has $r=1$, $\delta=0$ and begins by removing
the depth-one edge of arm two with endpoints $0$ and $8h+3$.
This is the source root edge $0$--$n$, since $n=b+2=8h+3$.

The complete finite complement of the $16$ rays is specified without
an arbitrary cutoff.  For each row $(r,\delta,b_0)$ in the threshold
table take
\begin{equation}\label{templates:min2-finite}
 \{(2,b,b+\delta):b\equiv r\pmod8,\ b<b_0,\
 b\ge5\text{ if }r\text{ is odd},\ b\ge6\text{ otherwise}\},
\end{equation}
and add $(2,3,9)$ and $(2,4,13)$.  The resulting set has exactly
$55$ profiles, all recorded in \texttt{F2.jsonl}.  Balanced odd
seeds have short leaf one; mixed even seeds used for resonance have
an upper leaf $K+e$ satisfying $2e\le b+4$.  These interfaces are
checked in addition to the alpha inventory.  The appendix's leaf
extension then reaches each resonance, and alpha hub attachment
handles the remaining larger long arm.  This specifies every
infinite ray and its exact finite remainder.

\subsubsection{The additional pinned-prefix condition for minimum three}
\label{templates:min3-protection}

The minimum-three proof first uses the same minimum-two templates,
but preserves the canonical first nine vertices of source arm one.
Every remove record has arm number one or two.  For a removed edge
on arm one with depth $ah+d$, add the condition $ah+d\ge10$.
All such slopes $a$ are positive in the actual records, so the bound
becomes
\[
 h'_0=\max\left(h_0,\max_{\text{arm-one cuts}}
                      \ceil{(10-d)/a}\right).
\]
Recomputing this expression on every record changes only
$(r,\delta)=(5,1)$: its threshold $b_0=21$ becomes $29$.
The endpoint-incidence identity adds no new incidence to a vertex
whose incident edges were frozen, so the protected prefix remains
the stated source prefix after the trade.

For $b\ge18$, the exact required pinned finite set is
\[
 \{(2,18,18)\}\ \cup\!!!
 \bigcup_{(r,\delta)}
 \{(2,b,b+\delta):18\le b<8h'_0(r,\delta)+r,\ b\equiv r\pmod8\}.
\]
It has $29$ profiles and is supplied completely in
\texttt{F3p.jsonl}.  Each has the maximum-pinned short word $(N,1)$;
every even mixed profile also has the upper-leaf bound $2e\le b+4$.
The even diagonal outside this finite set is handled by the separate
balanced constructions: $4\mid b$, $b\ge20$, or
$b\equiv2\pmod4$, $b\ge22$.  The latter has its first cut at
$b/2-1\ge10$.

For $3\le b\le17$, the exact direct minimum-three finite set is
\begin{align*}
 &\{(3,b,b+\delta):3\le b\le17,\ b+\delta\ge7,\\
 &\hspace{22mm}0\le\delta\le1\ (b\text{ odd}),\quad
                    0\le\delta\le2\ (b\text{ even})\}\\
 &\qquad\cup\ \{(3,b,4(1+\floor{b/2})+1):3\le b\le17\}.
\end{align*}
It has $44$ profiles, all in \texttt{F3s.jsonl}.  These sets are
derived from the thresholds and the attachment exception; they do
not introduce an empirical bound on the theorem.

\subsection{Two hundred sixteen minimum-four templates}
\label{templates:min4}

Now write $b=48h+r$, $0\le r<48$.  For odd $r$ take
$\delta=0,1,2,3$; for even $r$ take $\delta=0,1,2,3,4$.
There are $24\cdot4+24\cdot5=216$ records, one for each required
residue and near-band gap.

Put $n=b+\delta+4$ and $k=\floor{b/2}$.  Before raising upper
labels, the two-arm core is
\[
 (b,1,b-1,2),\qquad b-2-P_2(b-4).
\]
Raise every core label above $k$ by $n$.  The added arm is
\[
 \begin{cases}
 \operatorname{rev}(P_e(n))+k+1,&n\text{ even},\\
 n+k-\operatorname{rev}(P_e(n)),&n\text{ odd},
 \end{cases}
\qquad
 e=\begin{cases}\floor{n/2}-k-1,&n\text{ even},\\
 \floor{n/2}-k,&n\text{ odd}.
 \end{cases}
\]
For the required bands $e\in\{1,2,3\}$.  This completely defines
the source $\Sp{4,b-4,b+\delta+4}$, of total size
$N=2b+\delta+4$.  Its short word is $(N,1,N-1,2)$.
The kernel exception $b-4=9$ is excluded by the source domain;
the threshold or its literal finite complement handles that value.

Every removed edge in \texttt{S4.jsonl} has arm number one or two.
The endpoint-incidence identity then prevents any new incidence on
the frozen short arm.  The $216$ trades satisfy the exact color and
weighted-topology conditions with target lengths $(4,b,b+\delta)$.
The threshold table prints every $b_0$; the largest is $103$.
The thresholds arise from the explicit coordinate branches, not from
specializing the discovered words at a large parameter.

The required finite remainder is exactly
\begin{equation}\label{templates:min4-finite}
 \begin{aligned}
 \{(4,b,b+\delta):{}&b\ge9,\quad b\equiv r\pmod{48},\quad b<b_0(r,\delta),\\
                  &0\le\delta<4\text{ for odd }r,\quad
                   0\le\delta<5\text{ for even }r\}.
 \end{aligned}
\end{equation}
It has $139$ profiles, given in \texttt{F4.jsonl}.  Every word has
the pinned short arm $(N,1,N-1,2)$.  For odd $b$ at $\delta=0$,
the smaller upper long leaf $K+e$ satisfies $2e\le b+7$; for
even $b$ at $\delta=1$, the odd long leaf satisfies $2e\le b+8$.
These are the exact sufficient inequalities for the resonance
extensions.  The $20$ small auxiliary records for $b=4,\ldots,8$
are given in \texttt{F45.jsonl}; their individual pinned or alpha-only
roles are specified in the small-arm proof.

Finally replace a pinned four-word by
\[
 (N,1,N-1,2)\longmapsto(N+1,1,N,2,N-1).
\]
The label set gains exactly $N+1$; the old short palette
$[N-3,N]$ becomes $[N-3,N+1]$.  The cut and the other leaf labels
are unchanged.  Thus every pinned four-arm seed supplies its
minimum-five seed with the same resonance interface.  The $216$
ray certificates and their stated finite complements are therefore
the full finite input for both minimum-four and minimum-five bands.

\subsection{Equal-even repairs and the near EEO family}
\label{templates:equal-even}

Let $r$ be odd, $r\ge5$.  Start from the explicitly proved equal-odd
source $\Sp{(2r+1)^3}$, with arm order Walecki, A, C as in
Proposition~\ref{ok:equal-odd-three}.  Delete the last vertex of each arm;
their labels are $5r+3,5r+1,5r+2$.  Decrease the remaining labels
above $5r+3$ by three.  The shortened source has labels $[0,6r]$,
cut $3r$, three lengths $2r$, and defect
\[
 +\{4r+1,4r+2\}-\{4r,2r+4\}.
\]
Apply the lower and upper label cycles
\[
 (r-1,\ r+2,\ 3r-2,\ r+1,\ r),\qquad
 (4r-1,\ 4r+2,\ 4r+1,\ 4r+3),
\]
where each displayed value maps to the next, cyclically.  These
permutations stay within their respective classes.  The resulting
source has the exact defect
\begin{equation}\label{templates:equal-defect}
 +\{4r,2r+2,2r+1,2r+1\}
 -\{4,2r+5,4r-3,4r-2\}.
\end{equation}
Only the incident edges of the moved labels need to be inventoried
to obtain this identity.  On W their depths are $2r-2,2r-1,2r$;
on A they are $1,\ldots,6,2r-2,2r-1,2r$; on C they are
$1,2,2r$.  The explicit word coordinates prove that this is the
complete affected set.

For $r=4q+1$ and $r=4q+3$, $q\ge3$, the two entries of
\texttt{E.json} replace respectively $16$ and $15$ edges.  Their
full edge and source-position tables are printed below.  The signed
identity \eqref{templates:color-identity} cancels
\eqref{templates:equal-defect}; the weighted graph has three paths
of length $2r$.  The six small equal sources with lengths
$6,8,10,14,18,22$ are given explicitly in \texttt{Esmall.json}.
Balanced subdivision from these sources and the two rays gives
every even equal length at least six, as proved in the equal-triple
appendix.

The alternative $15$-edge table \texttt{Alt.json} uses the same
fixed-label source, but only $r=4q+1$, $q\ge3$.  It has the same
color and arm-length verification and an additional endpoint check:
the hub branch beginning at $H=5r+1$ ends at label $r$ and has
length $2r$.  The standard E recipe for $r=4q+3$ has this endpoint
property as well.  It is checked by following that particular
weighted hub branch and summing its exact affine weights.

\begin{lemma}\label{templates:reverse-extend}
Let a root-zero alpha spider have $N$ edges and an even arm whose
first and last labels are $H$ and $\ell$, with $H+\ell=N+1$.
Replacing that word by $(N+1,\operatorname{rev}(P))$ extends the
arm by one edge and preserves root-zero alpha labeling.
\end{lemma}
\begin{proof}
The new root color is $N+1$.  The next edge has color
$N+1-\ell=H$, replacing the former root color; all old internal
colors are unchanged.  The new label is the new maximum.  Reversing
an even word begins in the lower class, so prepending the new upper
vertex restores the alternating classes without changing the cut.
\end{proof}

Apply the lemma with $N=6r$, $H=5r+1$, and $\ell=r$ to obtain
$\Sp{2r,2r,2r+1}$.  The remaining parameters are fully specified:
the standard E recipe also works at the finite half-lengths
$r=7,11$.  For $r=3$, the equal-six arm
$(15,6,12,7,14,4)$ has endpoint sum $19=6r+1$, so the same
extension lemma applies with $H=15$, $\ell=4$.
Alt works directly at $r=9$; $a=8,10$ are literal near profiles.
If $a\equiv0\pmod4$
and $a\ge12$, apply Corollary~\ref{common:balanced-alpha-matching}
to the balanced equal source with lengths $a/2$. The unique-odd-arm
contree description in Corollary~\ref{common:unique-odd-contree}
identifies its matching lift with $\Sp{a,a,a+1}$.  Thus
\[
 a=8,\quad a=10,\quad a\ge12\ (4\mid a),\quad
 a\equiv6\pmod8,\quad a\ge18\ (a\equiv2\pmod8)
\]
are a disjoint exhaustive partition of the even $a\ge6$.
The literal targets at $a=6,8,10,14,18,22,26$ are also recorded
in \texttt{Nsmall.jsonl}.  They supplement the stated universal
proof rather than replacing it.

\subsection{Four equal odd arms}
\label{templates:four-odd}\label{construction:equal-odd-alpha4}

Write $s=4q+r$, $r\in\{0,1,2,3\}$, $q\ge6$, and let
$L=2s+1$.  Let $R_{s,i}(d)$, $i=0,1,2$, be the fully specified
equal-odd three-arm words of length $2s+1$ in the equal-triple
appendix, with $i=0$ their Walecki arm.  The unpermuted four-arm
source has the following complete coordinate rule for $d\ge1$:
\begin{equation}\label{templates:four-source}
 \begin{aligned}
 F_0(s,d)&=\begin{cases}
 8s+4-(d-1)/2,&d\text{ odd},\\ d/2,&d\text{ even},
 \end{cases}\\
 F_{i+1}(s,d)&=s+R_{s,i}(2s+2-d),\qquad i=0,1,2.
 \end{aligned}
\end{equation}
Set $F_i(s,0)=0$.  Interchange the two pairs of upper labels
$5s+1\leftrightarrow7s+1$ and $5s+2\leftrightarrow7s+2$.
These swaps preserve the label set, classes, and four source path
lengths.  Before the swaps the root-color defect is
$+\{6s+1,6s+2\}-\{4s+1,4s+2\}$.  Inventorying every edge incident
to the four moved labels gives the swapped seed defect
\begin{equation}\label{templates:four-defect}
 +\{4s-2,4s-1,4s,4s+1,6s+2\}
 -\{2s+1,2s+2,6s-2,6s-1,6s\}.
\end{equation}
The two labels $5s+1,5s+2$ are leaves and have one incident edge
each; $7s+1,7s+2$ have two.  Their complete local source coordinates
prove that no incident color is omitted in this calculation.
Explicitly, the six affected edges and colors are
\begin{center}\small
\begin{tabular}{llrr}
Old edge & New edge & Old color & New color\\\hline
$(3s-1,5s+1)$ & $(3s-1,7s+1)$ & $2s+2$ & $4s+2$\\
$(3s+1,5s+2)$ & $(3s+1,7s+2)$ & $2s+1$ & $4s+1$\\
$(s+1,7s+2)$ & $(s+1,5s+2)$ & $6s+1$ & $4s+1$\\
$(s+2,7s+2)$ & $(s+2,5s+2)$ & $6s$ & $4s$\\
$(s+2,7s+1)$ & $(s+2,5s+1)$ & $6s-1$ & $4s-1$\\
$(s+3,7s+1)$ & $(s+3,5s+1)$ & $6s-2$ & $4s-2$.
\end{tabular}\end{center}
The first two edges are the leaf edges of source arms two and three.
The last four occur on source arm one, at the reversals of depths
two through six of the original Walecki arm.  Thus the table follows
directly from \eqref{templates:four-source} and is exhaustive.  Adding
its signed changes to the pre-swap defect proves
\eqref{templates:four-defect}.

For each $r$, file \texttt{Qr.json} records the actual removed and
added edges and their positions in this swapped seed.  The four
trade sizes are
\begin{center}
\begin{tabular}{r|r|r|r}
$r$ & $L$ & Removed/added edges & Contracted vertices\\\hline
0 & $8q+1$ & $40/40$ &63\\
1 & $8q+3$ & $51/51$ &76\\
2 & $8q+5$ & $48/48$ &73\\
3 & $8q+7$ & $38/38$ &62.
\end{tabular}
\end{center}
For every entry, source-coordinate identities hold for all $q\ge6$,
the exact added/removed degree counters agree, and the affine color
calculation cancels \eqref{templates:four-defect}.  The full position
lists and both edge columns are printed below.  Weighted traversal
uses every contracted vertex and returns four paths of length
$8q+2r+1$, each as an affine identity.  The cut is $4s$ and the
vertex labels are $[0,8s+4]$.  Theorem~\ref{templates:soundness}
therefore proves root-zero alpha labeling for every odd $L\ge49$.

The remaining odd $L=7,9,\ldots,47$ have $21$ full literal records
in \texttt{Qsmall.jsonl}, checked by vertex, color, and alpha
inventories.  Together they give the equal-four input for every odd
$L\ge7$.  These finite records and the four rays are distinct
obligations; proving a ray does not certify its small-parameter
complement.

\subsection{Exact ray thresholds}
\label{templates:thresholds}
\begin{longtable}{rrrrr}
$b\bmod8$ & $\delta$ & $h_0$ & $b_0$ & Edges removed/added\\\hline\endhead
0 & 1 & 5 & 40 & 25\\
0 & 2 & 5 & 40 & 23\\
1 & 0 & 5 & 41 & 23\\
1 & 1 & 5 & 41 & 31\\
2 & 1 & 5 & 42 & 29\\
2 & 2 & 5 & 42 & 28\\
3 & 0 & 4 & 35 & 26\\
3 & 1 & 4 & 35 & 25\\
4 & 1 & 4 & 36 & 17\\
4 & 2 & 4 & 36 & 22\\
5 & 0 & 4 & 37 & 26\\
5 & 1 & 2 & 21 & 17\\
6 & 1 & 4 & 38 & 23\\
6 & 2 & 3 & 30 & 15\\
7 & 0 & 2 & 23 & 14\\
7 & 1 & 2 & 23 & 16\\
\end{longtable}
For the minimum-four rays, each entry is the inclusive threshold $b_0$; $h_0=(b_0-r)/48$. A dash denotes a band not required for that residue.
\begin{longtable}{r|rrrrr}
$r=b\bmod48$ & $\delta=0$ & $1$ & $2$ & $3$ & $4$\\\hline\endhead
0 & 48 & 48 & 48 & 48 & 48\\
1 & 49 & 49 & 49 & 49 & ---\\
2 & 50 & 50 & 50 & 50 & 50\\
3 & 51 & 51 & 51 & 51 & ---\\
4 & 52 & 52 & 52 & 52 & 52\\
5 & 53 & 53 & 53 & 53 & ---\\
6 & 54 & 54 & 54 & 102 & 54\\
7 & 55 & 55 & 103 & 55 & ---\\
8 & 56 & 56 & 56 & 56 & 56\\
9 & 57 & 57 & 57 & 57 & ---\\
10 & 58 & 58 & 58 & 58 & 58\\
11 & 59 & 59 & 59 & 59 & ---\\
12 & 60 & 60 & 60 & 60 & 60\\
13 & 61 & 61 & 61 & 61 & ---\\
14 & 62 & 62 & 62 & 62 & 62\\
15 & 63 & 63 & 63 & 63 & ---\\
16 & 64 & 64 & 64 & 64 & 64\\
17 & 65 & 65 & 65 & 65 & ---\\
18 & 66 & 66 & 66 & 66 & 66\\
19 & 67 & 67 & 67 & 67 & ---\\
20 & 68 & 68 & 68 & 68 & 68\\
21 & 69 & 69 & 69 & 69 & ---\\
22 & 70 & 70 & 70 & 70 & 70\\
23 & 71 & 71 & 71 & 71 & ---\\
24 & 72 & 72 & 72 & 72 & 72\\
25 & 73 & 73 & 73 & 73 & ---\\
26 & 74 & 74 & 74 & 74 & 74\\
27 & 75 & 75 & 75 & 75 & ---\\
28 & 76 & 76 & 76 & 76 & 76\\
29 & 77 & 77 & 77 & 77 & ---\\
30 & 78 & 78 & 78 & 78 & 78\\
31 & 79 & 79 & 79 & 79 & ---\\
32 & 80 & 80 & 80 & 80 & 80\\
33 & 81 & 81 & 81 & 81 & ---\\
34 & 82 & 82 & 82 & 82 & 82\\
35 & 35 & 83 & 83 & 83 & ---\\
36 & 36 & 84 & 36 & 36 & 36\\
37 & 85 & 85 & 37 & 85 & ---\\
38 & 38 & 38 & 38 & 38 & 38\\
39 & 39 & 39 & 39 & 39 & ---\\
40 & 40 & 40 & 40 & 40 & 40\\
41 & 41 & 41 & 41 & 41 & ---\\
42 & 90 & 90 & 90 & 90 & 90\\
43 & 43 & 43 & 91 & 91 & ---\\
44 & 44 & 44 & 44 & 44 & 44\\
45 & 93 & 93 & 93 & 93 & ---\\
46 & 46 & 46 & 46 & 46 & 94\\
47 & 47 & 47 & 47 & 95 & ---\\
\end{longtable}

\subsection{The complete equal-source permutation inventory}
\label{templates:equal-local-colors}
The following are all edges incident to a moved label in the two class cycles. Coordinates and colors are exact affine expressions in $q$, $q\ge3$. Subtracting the old column and adding the new column to the shortened-source defect gives the stated fixed-label seed defect.
\subsubsection[r=4q+1]{$r=4q+1$}
\begin{longtable}{rrrr}
Source arm & Edge depth & Old color & New color\\\hline\endhead
0 & $8q$ & $16q+7$ & $16q+4$\\
0 & $8q+1$ & $16q+6$ & $16q+3$\\
0 & $8q+2$ & $16q+5$ & $16q+6$\\
1 & $1$ & $16q+5$ & $16q+7$\\
1 & $2$ & $8q+4$ & $8q+6$\\
1 & $3$ & $8q+2$ & $8q+5$\\
1 & $4$ & $8q+1$ & $8q+4$\\
1 & $5$ & $8q+5$ & $8q+1$\\
1 & $6$ & $8q+7$ & $8q+3$\\
1 & $8q$ & $16q+1$ & $8q+3$\\
1 & $8q+1$ & $16q+2$ & $8q+4$\\
1 & $8q+2$ & $16q+3$ & $16q+4$\\
2 & $1$ & $16q+6$ & $16q+5$\\
2 & $2$ & $8q+3$ & $8q+2$\\
2 & $8q+2$ & $4$ & $8q+3$\\
\end{longtable}
\subsubsection[r=4q+3]{$r=4q+3$}
\begin{longtable}{rrrr}
Source arm & Edge depth & Old color & New color\\\hline\endhead
0 & $8q+4$ & $16q+15$ & $16q+12$\\
0 & $8q+5$ & $16q+14$ & $16q+11$\\
0 & $8q+6$ & $16q+13$ & $16q+14$\\
1 & $1$ & $16q+13$ & $16q+15$\\
1 & $2$ & $8q+8$ & $8q+10$\\
1 & $3$ & $8q+6$ & $8q+9$\\
1 & $4$ & $8q+5$ & $8q+8$\\
1 & $5$ & $8q+9$ & $8q+5$\\
1 & $6$ & $8q+11$ & $8q+7$\\
1 & $8q+4$ & $16q+9$ & $8q+7$\\
1 & $8q+5$ & $16q+10$ & $8q+8$\\
1 & $8q+6$ & $16q+11$ & $16q+12$\\
2 & $1$ & $16q+14$ & $16q+13$\\
2 & $2$ & $8q+7$ & $8q+6$\\
2 & $8q+6$ & $4$ & $8q+7$\\
\end{longtable}

\subsection{Complete equal-source and four-odd repair tables}
\label{templates:small-tables}
The vertex and depth conventions are those of the preceding schema. All edges in each remove column are replaced collectively by its add column.
\subsubsection[Equal-even source, r=4q+1, q >= 3]{Equal-even source, $r=4q+1$, $q\ge3$}
\label{templates:equal-r1}
\begingroup\small
\begin{longtable}{ll}
Remove & Add\\\hline\endhead
$(0,4r+3)$ & $(0,5r+1)$\\
$(\frac{1}{2}r-\frac{1}{2},\frac{11}{2}r+\frac{1}{2})$ & $(\frac{1}{2}r-\frac{1}{2},\frac{9}{2}r+\frac{3}{2})$\\
$(r-1,5r+1)$ & $(r-1,5r+2)$\\
$(r,5r)$ & $(r,4r+2)$\\
$(r+1,3r+2)$ & $(r+1,5r)$\\
$(r+2,5r+1)$ & $(r+2,5r-1)$\\
$(r+2,5r+2)$ & $(r+2,5r)$\\
$(\frac{3}{2}r-\frac{1}{2},\frac{9}{2}r+\frac{3}{2})$ & $(\frac{3}{2}r-\frac{1}{2},\frac{7}{2}r+\frac{5}{2})$\\
$(\frac{3}{2}r+\frac{1}{2},\frac{9}{2}r+\frac{3}{2})$ & $(\frac{3}{2}r+\frac{1}{2},\frac{11}{2}r+\frac{1}{2})$\\
$(2r-2,4r-1)$ & $(2r-2,4r+3)$\\
$(2r-1,4r+2)$ & $(2r-1,4r-1)$\\
$(2r,4r+2)$ & $(2r,5r+1)$\\
$(2r+1,4r+1)$ & $(2r+1,4r+2)$\\
$(\frac{5}{2}r-\frac{1}{2},\frac{7}{2}r+\frac{5}{2})$ & $(\frac{5}{2}r-\frac{1}{2},\frac{9}{2}r+\frac{3}{2})$\\
$(3r-2,5r-1)$ & $(3r-2,3r+2)$\\
$(3r-2,5r)$ & $(3r-2,4r+1)$\\
\end{longtable}
Rows form one collective edge trade; a row does not assert individual color preservation.
\begin{longtable}{rrr}
Vertex label & Source arm & Source depth\\\hline\endhead
$0$ & $\mathrm{hub}$ & $0$\\
$4r+3$ & 1 & $1$\\
$\frac{1}{2}r-\frac{1}{2}$ & 0 & $r-1$\\
$\frac{11}{2}r+\frac{1}{2}$ & 0 & $r$\\
$r-1$ & 0 & $2r$\\
$5r+1$ & 0 & $2r-1$\\
$r$ & 1 & $2r$\\
$5r$ & 1 & $2r-1$\\
$r+1$ & 2 & $2r$\\
$3r+2$ & 2 & $2r-1$\\
$r+2$ & 0 & $2r-2$\\
$5r+2$ & 0 & $2r-3$\\
$\frac{3}{2}r-\frac{1}{2}$ & 1 & $r+3$\\
$\frac{9}{2}r+\frac{3}{2}$ & 1 & $r+2$\\
$\frac{3}{2}r+\frac{1}{2}$ & 1 & $r+1$\\
$2r-2$ & 1 & $6$\\
$4r-1$ & 1 & $5$\\
$2r-1$ & 1 & $2$\\
$4r+2$ & 1 & $3$\\
$2r$ & 1 & $4$\\
$2r+1$ & 2 & $2$\\
$4r+1$ & 2 & $1$\\
$\frac{5}{2}r-\frac{1}{2}$ & 2 & $r-1$\\
$\frac{7}{2}r+\frac{5}{2}$ & 2 & $r$\\
$3r-2$ & 1 & $2r-2$\\
$5r-1$ & 1 & $2r-3$\\
\end{longtable}
\endgroup

\subsubsection[Equal-even source, r=4q+3, q >= 3]{Equal-even source, $r=4q+3$, $q\ge3$}
\label{templates:equal-r3}
\begingroup\small
\begin{longtable}{ll}
Remove & Add\\\hline\endhead
$(0,4r+3)$ & $(0,5r+1)$\\
$(\frac{1}{2}r-\frac{1}{2},\frac{11}{2}r+\frac{1}{2})$ & $(\frac{1}{2}r-\frac{1}{2},\frac{9}{2}r+\frac{3}{2})$\\
$(r-1,5r+1)$ & $(r-1,5r+2)$\\
$(r,5r)$ & $(r,4r+2)$\\
$(r+1,3r+2)$ & $(r+1,5r)$\\
$(r+2,5r+1)$ & $(r+2,5r-1)$\\
$(r+2,5r+2)$ & $(r+2,5r)$\\
$(\frac{3}{2}r-\frac{1}{2},\frac{9}{2}r+\frac{3}{2})$ & $(\frac{3}{2}r-\frac{1}{2},\frac{7}{2}r+\frac{3}{2})$\\
$(\frac{3}{2}r+\frac{1}{2},\frac{9}{2}r+\frac{1}{2})$ & $(\frac{3}{2}r+\frac{1}{2},\frac{11}{2}r+\frac{1}{2})$\\
$(2r-2,4r-1)$ & $(2r-2,4r+3)$\\
$(2r,4r+2)$ & $(2r,4r+1)$\\
$(2r+1,4r+1)$ & $(2r+1,5r+1)$\\
$(\frac{5}{2}r+\frac{1}{2},\frac{7}{2}r+\frac{3}{2})$ & $(\frac{5}{2}r+\frac{1}{2},\frac{9}{2}r+\frac{1}{2})$\\
$(3r-2,5r-1)$ & $(3r-2,3r+2)$\\
$(3r-2,5r)$ & $(3r-2,4r-1)$\\
\end{longtable}
Rows form one collective edge trade; a row does not assert individual color preservation.
\begin{longtable}{rrr}
Vertex label & Source arm & Source depth\\\hline\endhead
$0$ & $\mathrm{hub}$ & $0$\\
$4r+3$ & 1 & $1$\\
$\frac{1}{2}r-\frac{1}{2}$ & 0 & $r-1$\\
$\frac{11}{2}r+\frac{1}{2}$ & 0 & $r$\\
$r-1$ & 0 & $2r$\\
$5r+1$ & 0 & $2r-1$\\
$r$ & 1 & $2r$\\
$5r$ & 1 & $2r-1$\\
$r+1$ & 2 & $2r$\\
$3r+2$ & 2 & $2r-1$\\
$r+2$ & 0 & $2r-2$\\
$5r+2$ & 0 & $2r-3$\\
$\frac{3}{2}r-\frac{1}{2}$ & 1 & $r+3$\\
$\frac{9}{2}r+\frac{3}{2}$ & 1 & $r+2$\\
$\frac{3}{2}r+\frac{1}{2}$ & 1 & $r+1$\\
$\frac{9}{2}r+\frac{1}{2}$ & 1 & $r$\\
$2r-2$ & 1 & $6$\\
$4r-1$ & 1 & $5$\\
$2r$ & 1 & $4$\\
$4r+2$ & 1 & $3$\\
$2r+1$ & 2 & $2$\\
$4r+1$ & 2 & $1$\\
$\frac{5}{2}r+\frac{1}{2}$ & 2 & $r+1$\\
$\frac{7}{2}r+\frac{3}{2}$ & 2 & $r+2$\\
$3r-2$ & 1 & $2r-2$\\
$5r-1$ & 1 & $2r-3$\\
\end{longtable}
\endgroup

\subsubsection[Alternative complementary endpoints, r=4q+1, q >= 3]{Alternative complementary endpoints, $r=4q+1$, $q\ge3$}
\label{templates:alt}
\begingroup\small
\begin{longtable}{ll}
Remove & Add\\\hline\endhead
$(0,4r+3)$ & $(0,5r+1)$\\
$(\frac{1}{2}r-\frac{1}{2},\frac{11}{2}r+\frac{1}{2})$ & $(\frac{1}{2}r-\frac{1}{2},\frac{9}{2}r+\frac{3}{2})$\\
$(r-1,5r+1)$ & $(r-1,5r+2)$\\
$(r,5r)$ & $(r,4r+2)$\\
$(r+1,3r+2)$ & $(r+1,5r)$\\
$(r+2,5r+1)$ & $(r+2,5r-1)$\\
$(r+2,5r+2)$ & $(r+2,5r)$\\
$(\frac{3}{2}r-\frac{1}{2},\frac{9}{2}r+\frac{3}{2})$ & $(\frac{3}{2}r-\frac{1}{2},\frac{7}{2}r+\frac{5}{2})$\\
$(\frac{3}{2}r+\frac{1}{2},\frac{9}{2}r+\frac{1}{2})$ & $(\frac{3}{2}r+\frac{1}{2},\frac{11}{2}r+\frac{1}{2})$\\
$(2r-2,4r-1)$ & $(2r-2,4r+3)$\\
$(2r-1,4r+2)$ & $(2r-1,4r-1)$\\
$(2r+1,4r+1)$ & $(2r+1,5r+1)$\\
$(\frac{5}{2}r-\frac{1}{2},\frac{7}{2}r+\frac{5}{2})$ & $(\frac{5}{2}r-\frac{1}{2},\frac{9}{2}r+\frac{1}{2})$\\
$(3r-2,5r-1)$ & $(3r-2,3r+2)$\\
$(3r-2,5r)$ & $(3r-2,4r+1)$\\
\end{longtable}
Rows form one collective edge trade; a row does not assert individual color preservation.
\begin{longtable}{rrr}
Vertex label & Source arm & Source depth\\\hline\endhead
$0$ & $\mathrm{hub}$ & $0$\\
$\frac{1}{2}r-\frac{1}{2}$ & 0 & $r-1$\\
$r-1$ & 0 & $2r$\\
$r$ & 1 & $2r$\\
$r+1$ & 2 & $2r$\\
$r+2$ & 0 & $2r-2$\\
$\frac{3}{2}r-\frac{1}{2}$ & 1 & $r+3$\\
$\frac{3}{2}r+\frac{1}{2}$ & 1 & $r+1$\\
$2r-2$ & 1 & $6$\\
$2r-1$ & 1 & $2$\\
$2r+1$ & 2 & $2$\\
$\frac{5}{2}r-\frac{1}{2}$ & 2 & $r-1$\\
$3r-2$ & 1 & $2r-2$\\
$3r+2$ & 2 & $2r-1$\\
$\frac{7}{2}r+\frac{5}{2}$ & 2 & $r$\\
$4r-1$ & 1 & $5$\\
$4r+1$ & 2 & $1$\\
$4r+2$ & 1 & $3$\\
$4r+3$ & 1 & $1$\\
$\frac{9}{2}r+\frac{1}{2}$ & 1 & $r$\\
$\frac{9}{2}r+\frac{3}{2}$ & 1 & $r+2$\\
$5r-1$ & 1 & $2r-3$\\
$5r$ & 1 & $2r-1$\\
$5r+1$ & 0 & $2r-1$\\
$5r+2$ & 0 & $2r-3$\\
$\frac{11}{2}r+\frac{1}{2}$ & 0 & $r$\\
\end{longtable}
\endgroup

\subsubsection[Four equal odd arms, s=4q+0, q >= 6]{Four equal odd arms, $s=4q+0$, $q\ge6$}
\label{templates:four-r0}
\begingroup\small
\begin{longtable}{ll}
Remove & Add\\\hline\endhead
$(0,24q+2)$ & $(0,24q-1)$\\
$(4q+1,20q+2)$ & $(4q+1,24q+4)$\\
$(4q+2,20q+1)$ & $(4q+2,24q-2)$\\
$(4q+2,20q+2)$ & $(4q+2,28q+2)$\\
$(4q+3,20q+1)$ & $(4q+3,28q+1)$\\
$(6q,26q+3)$ & $(6q,26q+2)$\\
$(6q+1,26q+3)$ & $(6q+1,26q+1)$\\
$(6q+2,26q+2)$ & $(6q+2,26q)$\\
$(6q+3,26q+1)$ & $(6q+3,22q)$\\
$(6q+4,26q)$ & $(6q+4,22q-1)$\\
$(8q-1,24q+4)$ & $(8q-1,24q+3)$\\
$(8q,24q+3)$ & $(8q,24q+1)$\\
$(8q,24q+4)$ & $(8q,24q+2)$\\
$(8q+1,24q+1)$ & $(8q+1,20q+2)$\\
$(8q+2,24q-1)$ & $(8q+2,24q+2)$\\
$(8q+3,24q-2)$ & $(8q+3,20q-3)$\\
$(10q-2,22q+3)$ & $(10q-2,26q+3)$\\
$(10q,22q+1)$ & $(10q,26q+3)$\\
$(10q+1,22q)$ & $(10q+1,22q-1)$\\
$(10q+2,22q-1)$ & $(10q+2,22q-2)$\\
$(10q+2,22q)$ & $(10q+2,22q+1)$\\
$(10q+3,22q-1)$ & $(10q+3,22q)$\\
$(10q+4,22q-2)$ & $(10q+4,18q-1)$\\
$(12q-1,20q-1)$ & $(12q-1,20q+3)$\\
$(12q-1,28q+1)$ & $(12q-1,24q+4)$\\
$(12q,20q-1)$ & $(12q,20q+1)$\\
$(12q,20q+3)$ & $(12q,20q+2)$\\
$(12q+1,20q-3)$ & $(12q+1,20q)$\\
$(12q+1,28q+2)$ & $(12q+1,20q+1)$\\
$(12q+2,20q-3)$ & $(12q+2,20q-1)$\\
$(12q+3,20q)$ & $(12q+3,20q-1)$\\
$(14q,18q-1)$ & $(14q,22q+3)$\\
$(16q-5,16q+3)$ & $(16q-5,16q+4)$\\
$(16q-4,16q+3)$ & $(16q-4,16q+4)$\\
$(16q-3,16q+6)$ & $(16q-3,16q+3)$\\
$(16q-2,16q+4)$ & $(16q-2,16q+1)$\\
$(16q-2,24q+2)$ & $(16q-2,20q-3)$\\
$(16q-1,16q+1)$ & $(16q-1,16q+3)$\\
$(16q-1,16q+2)$ & $(16q-1,16q+6)$\\
$(16q,16q+4)$ & $(16q,16q+2)$\\
\end{longtable}
Rows form one collective edge trade; a row does not assert individual color preservation.
\begin{longtable}{rrr}
Vertex label & Source arm & Source depth\\\hline\endhead
$0$ & $\mathrm{hub}$ & $0$\\
$4q+1$ & 1 & $8q$\\
$4q+2$ & 1 & $8q-2$\\
$4q+3$ & 1 & $8q-4$\\
$6q$ & 1 & $4q+2$\\
$6q+1$ & 1 & $4q$\\
$6q+2$ & 1 & $4q-2$\\
$6q+3$ & 1 & $4q-4$\\
$6q+4$ & 1 & $4q-6$\\
$8q-1$ & 1 & $4$\\
$8q$ & 1 & $2$\\
$8q+1$ & 2 & $2$\\
$8q+2$ & 2 & $4$\\
$8q+3$ & 2 & $6$\\
$10q-2$ & 2 & $4q-4$\\
$10q$ & 2 & $4q$\\
$10q+1$ & 2 & $4q+2$\\
$10q+2$ & 2 & $4q+4$\\
$10q+3$ & 2 & $4q+6$\\
$10q+4$ & 2 & $4q+8$\\
$12q-1$ & 2 & $8q$\\
$12q$ & 2 & $8q-2$\\
$12q+1$ & 3 & $8q$\\
$12q+2$ & 3 & $8q-2$\\
$12q+3$ & 3 & $8q-4$\\
$14q$ & 3 & $4q+2$\\
$16q-5$ & 3 & $12$\\
$16q-4$ & 3 & $10$\\
$16q-3$ & 3 & $8$\\
$16q-2$ & 3 & $2$\\
$16q-1$ & 3 & $6$\\
$16q$ & 3 & $4$\\
$16q+1$ & 3 & $5$\\
$16q+2$ & 3 & $7$\\
$16q+3$ & 3 & $11$\\
$16q+4$ & 3 & $3$\\
$16q+6$ & 3 & $9$\\
$18q-1$ & 3 & $4q+3$\\
$20q-3$ & 3 & $8q-1$\\
$20q-1$ & 2 & $8q-1$\\
$20q$ & 3 & $8q-3$\\
$20q+1$ & 1 & $8q-3$\\
$20q+2$ & 1 & $8q-1$\\
$20q+3$ & 2 & $8q-3$\\
$22q-2$ & 2 & $4q+7$\\
$22q-1$ & 2 & $4q+5$\\
$22q$ & 2 & $4q+3$\\
$22q+1$ & 2 & $4q+1$\\
$22q+3$ & 2 & $4q-3$\\
$24q-2$ & 2 & $7$\\
$24q-1$ & 2 & $5$\\
$24q+1$ & 2 & $1$\\
$24q+2$ & 3 & $1$\\
$24q+3$ & 1 & $1$\\
$24q+4$ & 1 & $3$\\
$26q$ & 1 & $4q-5$\\
$26q+1$ & 1 & $4q-3$\\
$26q+2$ & 1 & $4q-1$\\
$26q+3$ & 1 & $4q+1$\\
$28q+1$ & 2 & $8q+1$\\
$28q+2$ & 3 & $8q+1$\\
\end{longtable}
\endgroup

\subsubsection[Four equal odd arms, s=4q+1, q >= 6]{Four equal odd arms, $s=4q+1$, $q\ge6$}
\label{templates:four-r1}
\begingroup\small
\begin{longtable}{ll}
Remove & Add\\\hline\endhead
$(0,24q+8)$ & $(0,24q+5)$\\
$(4q+2,20q+7)$ & $(4q+2,24q+8)$\\
$(4q+3,20q+6)$ & $(4q+3,24q+4)$\\
$(4q+3,20q+7)$ & $(4q+3,28q+9)$\\
$(4q+4,20q+6)$ & $(4q+4,28q+8)$\\
$(6q+3,26q+8)$ & $(6q+3,26q+6)$\\
$(6q+3,26q+9)$ & $(6q+3,26q+7)$\\
$(6q+4,26q+7)$ & $(6q+4,22q+5)$\\
$(6q+4,26q+8)$ & $(6q+4,26q+9)$\\
$(6q+5,26q+6)$ & $(6q+5,22q+4)$\\
$(8q+2,24q+9)$ & $(8q+2,24q+7)$\\
$(8q+3,24q+7)$ & $(8q+3,20q+8)$\\
$(8q+4,24q+5)$ & $(8q+4,24q+8)$\\
$(8q+5,24q+4)$ & $(8q+5,20q+2)$\\
$(10q+1,22q+9)$ & $(10q+1,26q+8)$\\
$(10q+2,22q+7)$ & $(10q+2,22q+6)$\\
$(10q+2,22q+8)$ & $(10q+2,26q+8)$\\
$(10q+3,22q+7)$ & $(10q+3,22q+9)$\\
$(10q+4,22q+5)$ & $(10q+4,18q+4)$\\
$(10q+4,22q+6)$ & $(10q+4,22q+3)$\\
$(10q+5,22q+4)$ & $(10q+5,22q+7)$\\
$(10q+6,22q+3)$ & $(10q+6,22q+7)$\\
$(12q+1,20q+8)$ & $(12q+1,24q+9)$\\
$(12q+2,20q+4)$ & $(12q+2,20q+7)$\\
$(12q+2,28q+8)$ & $(12q+2,20q+8)$\\
$(12q+3,20q+4)$ & $(12q+3,20q+6)$\\
$(12q+3,20q+8)$ & $(12q+3,20q+7)$\\
$(12q+4,20q+2)$ & $(12q+4,20q+5)$\\
$(12q+4,28q+9)$ & $(12q+4,20q+6)$\\
$(12q+5,20q+5)$ & $(12q+5,20q+4)$\\
$(12q+6,20q+5)$ & $(12q+6,20q+4)$\\
$(14q,18q+6)$ & $(14q,18q+7)$\\
$(14q+1,18q+6)$ & $(14q+1,22q+8)$\\
$(14q+2,18q+4)$ & $(14q+2,18q+6)$\\
$(14q+2,18q+9)$ & $(14q+2,18q+7)$\\
$(14q+3,18q+4)$ & $(14q+3,18q+5)$\\
$(14q+3,18q+7)$ & $(14q+3,18q+9)$\\
$(14q+4,18q+2)$ & $(14q+4,18q+3)$\\
$(14q+4,18q+7)$ & $(14q+4,18q+4)$\\
$(14q+5,18q+2)$ & $(14q+5,18q)$\\
$(14q+5,18q+5)$ & $(14q+5,18q+6)$\\
$(14q+6,18q)$ & $(14q+6,18q+2)$\\
$(14q+6,18q+5)$ & $(14q+6,18q+3)$\\
$(14q+7,18q+3)$ & $(14q+7,18q+5)$\\
$(14q+8,18q+3)$ & $(14q+8,18q+2)$\\
$(16q-2,16q+8)$ & $(16q-2,16q+9)$\\
$(16q,16q+11)$ & $(16q,16q+6)$\\
$(16q+1,16q+7)$ & $(16q+1,16q+11)$\\
$(16q+2,16q+6)$ & $(16q+2,16q+7)$\\
$(16q+2,24q+8)$ & $(16q+2,20q+5)$\\
$(16q+4,16q+9)$ & $(16q+4,16q+8)$\\
\end{longtable}
Rows form one collective edge trade; a row does not assert individual color preservation.
\begin{longtable}{rrr}
Vertex label & Source arm & Source depth\\\hline\endhead
$0$ & $\mathrm{hub}$ & $0$\\
$4q+2$ & 1 & $8q+2$\\
$4q+3$ & 1 & $8q$\\
$4q+4$ & 1 & $8q-2$\\
$6q+3$ & 1 & $4q$\\
$6q+4$ & 1 & $4q-2$\\
$6q+5$ & 1 & $4q-4$\\
$8q+2$ & 1 & $2$\\
$8q+3$ & 2 & $2$\\
$8q+4$ & 2 & $4$\\
$8q+5$ & 2 & $6$\\
$10q+1$ & 2 & $4q-2$\\
$10q+2$ & 2 & $4q$\\
$10q+3$ & 2 & $4q+2$\\
$10q+4$ & 2 & $4q+4$\\
$10q+5$ & 2 & $4q+6$\\
$10q+6$ & 2 & $4q+8$\\
$12q+1$ & 2 & $8q-2$\\
$12q+2$ & 2 & $8q+2$\\
$12q+3$ & 2 & $8q$\\
$12q+4$ & 3 & $8q+2$\\
$12q+5$ & 3 & $8q$\\
$12q+6$ & 3 & $8q-2$\\
$14q$ & 3 & $4q+10$\\
$14q+1$ & 3 & $4q+8$\\
$14q+2$ & 3 & $4q+6$\\
$14q+3$ & 3 & $4q+4$\\
$14q+4$ & 3 & $4q+2$\\
$14q+5$ & 3 & $4q$\\
$14q+6$ & 3 & $4q-2$\\
$14q+7$ & 3 & $4q-4$\\
$14q+8$ & 3 & $4q-6$\\
$16q-2$ & 3 & $14$\\
$16q$ & 3 & $10$\\
$16q+1$ & 3 & $8$\\
$16q+2$ & 3 & $2$\\
$16q+4$ & 3 & $6$\\
$16q+6$ & 3 & $3$\\
$16q+7$ & 3 & $9$\\
$16q+8$ & 3 & $13$\\
$16q+9$ & 3 & $7$\\
$16q+11$ & 3 & $11$\\
$18q$ & 3 & $4q-3$\\
$18q+2$ & 3 & $4q+1$\\
$18q+3$ & 3 & $4q-5$\\
$18q+4$ & 3 & $4q+5$\\
$18q+5$ & 3 & $4q-1$\\
$18q+6$ & 3 & $4q+9$\\
$18q+7$ & 3 & $4q+3$\\
$18q+9$ & 3 & $4q+7$\\
$20q+2$ & 3 & $8q+1$\\
$20q+4$ & 2 & $8q+1$\\
$20q+5$ & 3 & $8q-1$\\
$20q+6$ & 1 & $8q-1$\\
$20q+7$ & 1 & $8q+1$\\
$20q+8$ & 2 & $8q-1$\\
$22q+3$ & 2 & $4q+9$\\
$22q+4$ & 2 & $4q+7$\\
$22q+5$ & 2 & $4q+5$\\
$22q+6$ & 2 & $4q+3$\\
$22q+7$ & 2 & $4q+1$\\
$22q+8$ & 2 & $4q-1$\\
$22q+9$ & 2 & $4q-3$\\
$24q+4$ & 2 & $7$\\
$24q+5$ & 2 & $5$\\
$24q+7$ & 2 & $1$\\
$24q+8$ & 3 & $1$\\
$24q+9$ & 1 & $1$\\
$26q+6$ & 1 & $4q-5$\\
$26q+7$ & 1 & $4q-3$\\
$26q+8$ & 1 & $4q-1$\\
$26q+9$ & 1 & $4q+1$\\
$28q+8$ & 2 & $8q+3$\\
$28q+9$ & 3 & $8q+3$\\
\end{longtable}
\endgroup

\subsubsection[Four equal odd arms, s=4q+2, q >= 6]{Four equal odd arms, $s=4q+2$, $q\ge6$}
\label{templates:four-r2}
\begingroup\small
\begin{longtable}{ll}
Remove & Add\\\hline\endhead
$(0,24q+13)$ & $(0,24q+10)$\\
$(4q+3,20q+12)$ & $(4q+3,24q+16)$\\
$(4q+3,28q+17)$ & $(4q+3,28q+16)$\\
$(4q+4,20q+11)$ & $(4q+4,24q+11)$\\
$(4q+4,20q+12)$ & $(4q+4,28q+15)$\\
$(4q+5,20q+11)$ & $(4q+5,28q+17)$\\
$(6q+3,26q+16)$ & $(6q+3,26q+14)$\\
$(6q+4,26q+15)$ & $(6q+4,26q+14)$\\
$(6q+5,26q+14)$ & $(6q+5,22q+10)$\\
$(6q+5,26q+15)$ & $(6q+5,26q+13)$\\
$(6q+6,26q+13)$ & $(6q+6,22q+9)$\\
$(6q+6,26q+14)$ & $(6q+6,26q+15)$\\
$(8q+3,24q+16)$ & $(8q+3,24q+13)$\\
$(8q+4,24q+16)$ & $(8q+4,24q+13)$\\
$(8q+5,24q+13)$ & $(8q+5,20q+12)$\\
$(8q+6,24q+11)$ & $(8q+6,24q+14)$\\
$(8q+7,24q+10)$ & $(8q+7,20q+7)$\\
$(10q+3,22q+14)$ & $(10q+3,26q+15)$\\
$(10q+3,22q+15)$ & $(10q+3,26q+16)$\\
$(10q+4,22q+13)$ & $(10q+4,22q+15)$\\
$(10q+5,22q+12)$ & $(10q+5,22q+14)$\\
$(10q+7,22q+10)$ & $(10q+7,18q+9)$\\
$(10q+9,22q+9)$ & $(10q+9,22q+12)$\\
$(12q+4,20q+13)$ & $(12q+4,24q+16)$\\
$(12q+5,20q+9)$ & $(12q+5,20q+12)$\\
$(12q+5,28q+15)$ & $(12q+5,20q+13)$\\
$(12q+6,20q+9)$ & $(12q+6,20q+10)$\\
$(12q+6,20q+13)$ & $(12q+6,20q+11)$\\
$(12q+7,20q+7)$ & $(12q+7,20q+10)$\\
$(12q+7,28q+16)$ & $(12q+7,20q+13)$\\
$(12q+8,20q+10)$ & $(12q+8,20q+9)$\\
$(12q+9,20q+10)$ & $(12q+9,20q+9)$\\
$(14q+4,18q+11)$ & $(14q+4,22q+13)$\\
$(14q+5,18q+9)$ & $(14q+5,18q+12)$\\
$(14q+6,18q+9)$ & $(14q+6,18q+7)$\\
$(14q+7,18q+7)$ & $(14q+7,18q+10)$\\
$(14q+7,18q+12)$ & $(14q+7,18q+11)$\\
$(14q+8,18q+7)$ & $(14q+8,18q+8)$\\
$(14q+9,18q+10)$ & $(14q+9,18q+8)$\\
$(14q+10,18q+8)$ & $(14q+10,18q+7)$\\
$(14q+11,18q+8)$ & $(14q+11,18q+9)$\\
$(16q+3,16q+11)$ & $(16q+3,16q+12)$\\
$(16q+4,16q+11)$ & $(16q+4,16q+12)$\\
$(16q+5,16q+14)$ & $(16q+5,16q+11)$\\
$(16q+6,16q+12)$ & $(16q+6,16q+10)$\\
$(16q+6,24q+14)$ & $(16q+6,20q+11)$\\
$(16q+7,16q+10)$ & $(16q+7,16q+14)$\\
$(16q+8,16q+12)$ & $(16q+8,16q+11)$\\
\end{longtable}
Rows form one collective edge trade; a row does not assert individual color preservation.
\begin{longtable}{rrr}
Vertex label & Source arm & Source depth\\\hline\endhead
$0$ & $\mathrm{hub}$ & $0$\\
$4q+3$ & 1 & $8q+4$\\
$4q+4$ & 1 & $8q+2$\\
$4q+5$ & 1 & $8q$\\
$6q+3$ & 1 & $4q+4$\\
$6q+4$ & 1 & $4q+2$\\
$6q+5$ & 1 & $4q$\\
$6q+6$ & 1 & $4q-2$\\
$8q+3$ & 1 & $4$\\
$8q+4$ & 1 & $2$\\
$8q+5$ & 2 & $2$\\
$8q+6$ & 2 & $4$\\
$8q+7$ & 2 & $6$\\
$10q+3$ & 2 & $4q-2$\\
$10q+4$ & 2 & $4q$\\
$10q+5$ & 2 & $4q+2$\\
$10q+7$ & 2 & $4q+6$\\
$10q+9$ & 2 & $4q+10$\\
$12q+4$ & 2 & $8q$\\
$12q+5$ & 2 & $8q+4$\\
$12q+6$ & 2 & $8q+2$\\
$12q+7$ & 3 & $8q+4$\\
$12q+8$ & 3 & $8q+2$\\
$12q+9$ & 3 & $8q$\\
$14q+4$ & 3 & $4q+10$\\
$14q+5$ & 3 & $4q+8$\\
$14q+6$ & 3 & $4q+6$\\
$14q+7$ & 3 & $4q+4$\\
$14q+8$ & 3 & $4q+2$\\
$14q+9$ & 3 & $4q$\\
$14q+10$ & 3 & $4q-2$\\
$14q+11$ & 3 & $4q-4$\\
$16q+3$ & 3 & $12$\\
$16q+4$ & 3 & $10$\\
$16q+5$ & 3 & $8$\\
$16q+6$ & 3 & $2$\\
$16q+7$ & 3 & $6$\\
$16q+8$ & 3 & $4$\\
$16q+10$ & 3 & $7$\\
$16q+11$ & 3 & $11$\\
$16q+12$ & 3 & $3$\\
$16q+14$ & 3 & $9$\\
$18q+7$ & 3 & $4q+3$\\
$18q+8$ & 3 & $4q-3$\\
$18q+9$ & 3 & $4q+7$\\
$18q+10$ & 3 & $4q+1$\\
$18q+11$ & 3 & $4q+11$\\
$18q+12$ & 3 & $4q+5$\\
$20q+7$ & 3 & $8q+3$\\
$20q+9$ & 2 & $8q+3$\\
$20q+10$ & 3 & $8q+1$\\
$20q+11$ & 1 & $8q+1$\\
$20q+12$ & 1 & $8q+3$\\
$20q+13$ & 2 & $8q+1$\\
$22q+9$ & 2 & $4q+9$\\
$22q+10$ & 2 & $4q+7$\\
$22q+12$ & 2 & $4q+3$\\
$22q+13$ & 2 & $4q+1$\\
$22q+14$ & 2 & $4q-1$\\
$22q+15$ & 2 & $4q-3$\\
$24q+10$ & 2 & $7$\\
$24q+11$ & 2 & $5$\\
$24q+13$ & 2 & $1$\\
$24q+14$ & 3 & $1$\\
$24q+16$ & 1 & $3$\\
$26q+13$ & 1 & $4q-3$\\
$26q+14$ & 1 & $4q-1$\\
$26q+15$ & 1 & $4q+1$\\
$26q+16$ & 1 & $4q+3$\\
$28q+15$ & 2 & $8q+5$\\
$28q+16$ & 3 & $8q+5$\\
$28q+17$ & 1 & $8q+5$\\
\end{longtable}
\endgroup

\subsubsection[Four equal odd arms, s=4q+3, q >= 6]{Four equal odd arms, $s=4q+3$, $q\ge6$}
\label{templates:four-r3}
\begingroup\small
\begin{longtable}{ll}
Remove & Add\\\hline\endhead
$(0,24q+19)$ & $(0,24q+16)$\\
$(4q+4,20q+17)$ & $(4q+4,24q+22)$\\
$(4q+4,28q+24)$ & $(4q+4,28q+22)$\\
$(4q+5,20q+16)$ & $(4q+5,24q+17)$\\
$(4q+5,20q+17)$ & $(4q+5,28q+24)$\\
$(4q+6,20q+16)$ & $(4q+6,28q+23)$\\
$(6q+5,26q+23)$ & $(6q+5,22q+19)$\\
$(6q+8,26q+20)$ & $(6q+8,22q+15)$\\
$(8q+5,24q+22)$ & $(8q+5,24q+20)$\\
$(8q+6,24q+21)$ & $(8q+6,24q+19)$\\
$(8q+8,24q+17)$ & $(8q+8,20q+16)$\\
$(8q+9,24q+16)$ & $(8q+9,20q+12)$\\
$(10q+6,22q+20)$ & $(10q+6,26q+23)$\\
$(10q+7,22q+19)$ & $(10q+7,22q+17)$\\
$(10q+8,22q+18)$ & $(10q+8,22q+20)$\\
$(10q+9,22q+17)$ & $(10q+9,22q+15)$\\
$(10q+10,22q+15)$ & $(10q+10,18q+14)$\\
$(10q+10,22q+16)$ & $(10q+10,22q+14)$\\
$(10q+11,22q+14)$ & $(10q+11,22q+16)$\\
$(10q+11,22q+15)$ & $(10q+11,26q+20)$\\
$(12q+7,20q+18)$ & $(12q+7,24q+21)$\\
$(12q+8,20q+14)$ & $(12q+8,20q+16)$\\
$(12q+8,28q+22)$ & $(12q+8,20q+18)$\\
$(12q+9,20q+14)$ & $(12q+9,20q+15)$\\
$(12q+10,20q+12)$ & $(12q+10,20q+15)$\\
$(12q+10,28q+23)$ & $(12q+10,20q+17)$\\
$(12q+11,20q+15)$ & $(12q+11,20q+14)$\\
$(12q+12,20q+15)$ & $(12q+12,20q+14)$\\
$(14q+7,18q+16)$ & $(14q+7,22q+18)$\\
$(14q+8,18q+14)$ & $(14q+8,18q+17)$\\
$(14q+9,18q+14)$ & $(14q+9,18q+12)$\\
$(14q+10,18q+12)$ & $(14q+10,18q+15)$\\
$(14q+10,18q+17)$ & $(14q+10,18q+16)$\\
$(14q+11,18q+12)$ & $(14q+11,18q+13)$\\
$(14q+12,18q+15)$ & $(14q+12,18q+13)$\\
$(14q+13,18q+13)$ & $(14q+13,18q+12)$\\
$(14q+14,18q+13)$ & $(14q+14,18q+14)$\\
$(16q+10,24q+20)$ & $(16q+10,20q+17)$\\
\end{longtable}
Rows form one collective edge trade; a row does not assert individual color preservation.
\begin{longtable}{rrr}
Vertex label & Source arm & Source depth\\\hline\endhead
$0$ & $\mathrm{hub}$ & $0$\\
$4q+4$ & 1 & $8q+6$\\
$4q+5$ & 1 & $8q+4$\\
$4q+6$ & 1 & $8q+2$\\
$6q+5$ & 1 & $4q+4$\\
$6q+8$ & 1 & $4q-2$\\
$8q+5$ & 1 & $4$\\
$8q+6$ & 1 & $2$\\
$8q+8$ & 2 & $4$\\
$8q+9$ & 2 & $6$\\
$10q+6$ & 2 & $4q$\\
$10q+7$ & 2 & $4q+2$\\
$10q+8$ & 2 & $4q+4$\\
$10q+9$ & 2 & $4q+6$\\
$10q+10$ & 2 & $4q+8$\\
$10q+11$ & 2 & $4q+10$\\
$12q+7$ & 2 & $8q+2$\\
$12q+8$ & 2 & $8q+6$\\
$12q+9$ & 2 & $8q+4$\\
$12q+10$ & 3 & $8q+6$\\
$12q+11$ & 3 & $8q+4$\\
$12q+12$ & 3 & $8q+2$\\
$14q+7$ & 3 & $4q+12$\\
$14q+8$ & 3 & $4q+10$\\
$14q+9$ & 3 & $4q+8$\\
$14q+10$ & 3 & $4q+6$\\
$14q+11$ & 3 & $4q+4$\\
$14q+12$ & 3 & $4q+2$\\
$14q+13$ & 3 & $4q$\\
$14q+14$ & 3 & $4q-2$\\
$16q+10$ & 3 & $2$\\
$18q+12$ & 3 & $4q+5$\\
$18q+13$ & 3 & $4q-1$\\
$18q+14$ & 3 & $4q+9$\\
$18q+15$ & 3 & $4q+3$\\
$18q+16$ & 3 & $4q+13$\\
$18q+17$ & 3 & $4q+7$\\
$20q+12$ & 3 & $8q+5$\\
$20q+14$ & 2 & $8q+5$\\
$20q+15$ & 3 & $8q+3$\\
$20q+16$ & 1 & $8q+3$\\
$20q+17$ & 1 & $8q+5$\\
$20q+18$ & 2 & $8q+3$\\
$22q+14$ & 2 & $4q+11$\\
$22q+15$ & 2 & $4q+9$\\
$22q+16$ & 2 & $4q+7$\\
$22q+17$ & 2 & $4q+5$\\
$22q+18$ & 2 & $4q+3$\\
$22q+19$ & 2 & $4q+1$\\
$22q+20$ & 2 & $4q-1$\\
$24q+16$ & 2 & $7$\\
$24q+17$ & 2 & $5$\\
$24q+19$ & 2 & $1$\\
$24q+20$ & 3 & $1$\\
$24q+21$ & 1 & $1$\\
$24q+22$ & 1 & $3$\\
$26q+20$ & 1 & $4q-1$\\
$26q+23$ & 1 & $4q+5$\\
$28q+22$ & 2 & $8q+7$\\
$28q+23$ & 3 & $8q+7$\\
$28q+24$ & 1 & $8q+7$\\
\end{longtable}
\endgroup


\subsection{Complete finite input ledger}
\label{templates:ledger}
The aliases below are the files in \texttt{data/affine/} of the supporting repository.
\begingroup\footnotesize\begin{longtable}{lr}
Input & SHA-256 (first and last halves)\\\hline\endhead
\texttt{S2o.jsonl} & \texttt{9f1f9481846b81323feae9828c6fcb71}\\
 & \texttt{4ba7f2be21cecbd9d09c385930005b2f}\\[3pt]
\texttt{S2e.jsonl} & \texttt{841cc9bfe43dd853a4c1185176a729f2}\\
 & \texttt{c30ca678ceb7abe6abd2bfda48f69ac3}\\[3pt]
\texttt{S4.jsonl} & \texttt{60b39f453abdc6fa130429c1df76457b}\\
 & \texttt{ca0a6294f9f93ea3bb43bfe00de3d229}\\[3pt]
\texttt{E.json} & \texttt{0eefe79d40a9df8aae39132d4b782681}\\
 & \texttt{bba19bf801c37e9bfd7bc8fd18dc37dc}\\[3pt]
\texttt{Alt.json} & \texttt{655183b5952106706804d3010b3be969}\\
 & \texttt{010cd5ff8a18d94871c55d56299804a7}\\[3pt]
\texttt{Q0.json} & \texttt{07b2679b47f3246af2d8776674679fed}\\
 & \texttt{d4535a262ee4ea2787c4c58afb7ebc9a}\\[3pt]
\texttt{Q1.json} & \texttt{d55627eca8bf2b27ff66657bb6e282a9}\\
 & \texttt{c9aae4ee841f344854883d002c675131}\\[3pt]
\texttt{Q2.json} & \texttt{0f165f3cea4ed0d71cf105cadffe8f2c}\\
 & \texttt{2ebb031e871ffecd6531db9490806317}\\[3pt]
\texttt{Q3.json} & \texttt{99cdc386444ce9f992a7266cd5e59642}\\
 & \texttt{34c5045cc6c7277ceb2f04154ae13582}\\[3pt]
\texttt{P.json} & \texttt{9b5d5d28c24bc60a52a63039c17c6786}\\
 & \texttt{bdcb0e13c135eac070e65a0aabccd859}\\[3pt]
\texttt{F2.jsonl} & \texttt{0e67af0313b85faac029eb0aca15c6a6}\\
 & \texttt{9b2c8b65a3df07eeac8927a6663aa909}\\[3pt]
\texttt{F3s.jsonl} & \texttt{28520669a674bf72dc9c2cb854ac4fe6}\\
 & \texttt{abed99be794d026b441453269db400ad}\\[3pt]
\texttt{F3p.jsonl} & \texttt{114e36ecda7dcba694628338bdde0aa3}\\
 & \texttt{cbf11eb4c002d708780303b39712eba7}\\[3pt]
\texttt{F4.jsonl} & \texttt{3df86054602a198e57ebe553fff61de7}\\
 & \texttt{26b77153e2bb8d08205170e430f59532}\\[3pt]
\texttt{F45.jsonl} & \texttt{0561aacc63e517c97f9c76e9f54f5394}\\
 & \texttt{90c6032551a66c89eeff40cadc4f65c8}\\[3pt]
\texttt{Nsmall.jsonl} & \texttt{600fd54ba59790f4315dd2c97f024e4e}\\
 & \texttt{c5d9c3a60520b7bd118df184fc4e9190}\\[3pt]
\texttt{Qsmall.jsonl} & \texttt{383ee17c3b04c362f435cc80aae39541}\\
 & \texttt{3d4483824d693bac23873d4c4fcb975b}\\[3pt]
\texttt{Esmall.json} & \texttt{8c9d47773f4c44b29b1a64aa8e95814c}\\
 & \texttt{122f6341838d040bd92bfa3c258767ef}\\[3pt]
\end{longtable}\endgroup

